The question concerns the explanatory status of constants. When a
theory receives a parameter determined experimentally, its equations
establish the consequences that follow from that value. A structural
derivation changes this relationship: the value becomes one of the
consequences of the construction. The question of the strength of a
coupling or of a characteristic ratio of matter then becomes
the question of the mathematical relations that determine them.
The transition from parameter to consequence constitutes the guiding argument
of this article.

Numerical determination acquires, from this perspective, a constitutive
content. The equation brings together operations whose provenance must be made
explicit; its evaluation obtains the value fixed by those operations. Once
the corresponding conditions of existence, uniqueness, and compatibility
have been established, the necessity of the result belongs to the proof.
Explaining why a constant takes a value means reconstructing that
necessity within the domain of the theory. Decimal precision shows
a quantitative consequence of the structure, while the derivation
sets out the reason why precisely that consequence is obtained.

The scope of this question increases when several determinations
share a common origin. The thesis proposed by the authors links
constants to the organization that produces their reciprocal relations.
The common architecture must explain both the individuality of each
coordinate and the dependencies linking it to the others.
The order of the constructions therefore acquires demonstrative value:
it makes it possible to recognize which data are generated, which operations use them, and
which relations survive successive realizations. The diversity
of values is studied as an expression of an articulated formal unity.

In this organization, metrology performs a subsequent and
essential function. The construction fixes its results before comparing them with
observed quantities; experience examines their physical
correspondence and delimits the scope of empirical agreement. Mathematical necessity is established
with respect to the declared premises and operations, and natural
realization is examined through the identification of observables and their
independent testing. Both tasks contribute to a complete
explanation: the first determines the result within the corpus; the second
studies its relevance to describing the quantitative constitution of the
world.

This perspective makes it possible to situate the research historically. The
precedents considered below specify different ways of
relating laws, parameters, and observation. The present work adopts
as its central problem the internal determination of values and of the
structure linking them. Its arithmetic, geometric, and
algebraic developments receive their argumentative unity from that question. The
ontological aspiration lies in explaining how apparently independent
physical properties may correspond to necessary relations of a single
mathematical organization.
